\documentclass[12pt, letterpaper]{article}

% ── Paquetes ──────────────────────────────────────────────────────────────
\usepackage[utf8]{inputenc}
\usepackage[T1]{fontenc}
\usepackage[spanish]{babel}
\usepackage{geometry}
\usepackage{amsmath, amssymb, amsthm}
\usepackage{booktabs}
\usepackage{array}
\usepackage{microtype}
\usepackage{setspace}
\usepackage{parskip}
\usepackage{titlesec}
\usepackage{fancyhdr}
\usepackage{enumitem}
\usepackage{bm}
\usepackage{mathtools}
\usepackage{tcolorbox}
\usepackage{hyperref}

\tcbuselibrary{skins, breakable}

% ── Geometría ─────────────────────────────────────────────────────────────
\geometry{top=2.8cm, bottom=3.0cm, left=3.2cm, right=3.2cm}
\setstretch{1.20}

% ── Encabezado ────────────────────────────────────────────────────────────
\pagestyle{fancy}
\fancyhf{}
\fancyhead[C]{\small\textsc{Econometría · Gujarati — Ejercicio 10.4}}
\fancyfoot[C]{\small\thepage}
\renewcommand{\headrulewidth}{0.4pt}
\renewcommand{\footrulewidth}{0pt}

% ── Estilos de sección estilo NYT ─────────────────────────────────────────
\titleformat{\section}
  {\large\bfseries\scshape}{\thesection.}{0.6em}{}
  [\vspace{-4pt}\rule{\linewidth}{0.4pt}]

\titleformat{\subsection}
  {\normalsize\bfseries}{\thesubsection.}{0.5em}{}

\titleformat{\subsubsection}
  {\normalsize\itshape}{}{0em}{}

% ── Entornos ──────────────────────────────────────────────────────────────
\newtheorem{prop}{Proposición}
\newtheorem{lema}{Lema}
\theoremstyle{remark}
\newtheorem{obs}{Observación}

% Caja de resultado clave
\newtcolorbox{clave}{
  colback=white, colframe=black,
  boxrule=0.5pt, arc=3pt,
  left=8pt, right=8pt, top=6pt, bottom=6pt,
  fonttitle=\bfseries\small,
  breakable
}

% ── Comandos ──────────────────────────────────────────────────────────────
\newcommand{\X}{\mathbf{X}}
\newcommand{\y}{\mathbf{y}}
\newcommand{\bhat}{\hat{\boldsymbol{\beta}}}
\newcommand{\lam}{\lambda}
\newcommand{\Cov}{\operatorname{Cov}}
\newcommand{\Var}{\operatorname{Var}}
\newcommand{\Corr}{\operatorname{Corr}}
\newcommand{\plim}{\operatorname{plim}}
\newcommand{\tr}{\operatorname{tr}}
\newcommand{\rp}[2]{r_{#1 \cdot #2}}   % correlación parcial
\newcommand{\Rd}[2]{R^2_{#1 \cdot #2}} % R² parcial

% ─────────────────────────────────────────────────────────────────────────
\begin{document}

% ── Portada ───────────────────────────────────────────────────────────────
\begin{center}
  {\LARGE\bfseries Multicolinealidad Perfecta:}\\[6pt]
  {\LARGE\bfseries Correlaciones Parciales y Coeficientes}\\
  {\LARGE\bfseries de Determinación bajo Dependencia Lineal Exacta}\\[12pt]
  {\large\itshape Derivación completa desde primeros principios}\\[14pt]
  \rule{0.55\linewidth}{0.6pt}\\[8pt]
  {\normalsize Ejercicio~10.4 — \textit{Econometría}, Gujarati \& Porter, 5.ª ed.}\\[4pt]
  {\small\textsc{Con álgebra matricial, geometría de correlación y todas las demostraciones}}
\end{center}

\vspace{10pt}

\noindent\textbf{Planteamiento del problema.}\quad
Se afirma que la relación
\begin{equation}
  \lambda_1 X_{1i} + \lambda_2 X_{2i} + \lambda_3 X_{3i} = 0
  \qquad \text{para todo } i = 1, \ldots, n
  \label{eq:relacion}
\end{equation}
se mantiene con $\lambda_1, \lambda_2, \lambda_3$ no todos simultáneamente cero.
Se pide:
\begin{enumerate}[label=(\alph*), itemsep=3pt]
  \item Estimar las correlaciones parciales $\rp{12}{3}$, $\rp{13}{2}$ y $\rp{23}{1}$.
  \item Encontrar los coeficientes de determinación $\Rd{1}{23}$, $\Rd{2}{13}$ y $\Rd{3}{12}$.
  \item Determinar el grado de multicolinealidad en esta situación.
\end{enumerate}

% ══════════════════════════════════════════════════════════════════════════
\section{Fundamentos y definiciones}

\subsection{Correlación simple entre dos variables}

Dadas $n$ observaciones de dos variables $X_j$ y $X_k$, la \textbf{correlación
de Pearson} es:

\begin{equation}
  r_{jk}
    = \frac{\displaystyle\sum_{i=1}^{n}(X_{ji}-\bar{X}_j)(X_{ki}-\bar{X}_k)}
           {\sqrt{\displaystyle\sum_{i=1}^{n}(X_{ji}-\bar{X}_j)^2}
            \sqrt{\displaystyle\sum_{i=1}^{n}(X_{ki}-\bar{X}_k)^2}}
    = \frac{\sum x_{ji}\,x_{ki}}{\sqrt{\sum x_{ji}^2\,\sum x_{ki}^2}}
  \label{eq:pearson}
\end{equation}

\noindent donde $x_{ji} = X_{ji} - \bar{X}_j$ denota la desviación respecto a la
media. Por la desigualdad de Cauchy-Schwarz, $-1 \le r_{jk} \le 1$, con
$|r_{jk}| = 1$ si y sólo si existe una relación lineal exacta entre $X_j$ y
$X_k$.

\subsection{Correlación parcial: concepto y fórmula}

La \textbf{correlación parcial} $\rp{jk}{\ell}$ mide la asociación lineal entre
$X_j$ y $X_k$ una vez que se ha eliminado la influencia de una tercera variable
$X_\ell$. Formalmente, sea:

\begin{align}
  \tilde{X}_{ji} &= X_{ji} - \hat{X}_{ji}^{(\ell)}
  \label{eq:resid_j}\\
  \tilde{X}_{ki} &= X_{ki} - \hat{X}_{ki}^{(\ell)}
  \label{eq:resid_k}
\end{align}

\noindent donde $\hat{X}_{ji}^{(\ell)}$ y $\hat{X}_{ki}^{(\ell)}$ son los
valores ajustados de las regresiones de $X_j$ y $X_k$ sobre $X_\ell$,
respectivamente. Entonces:

\begin{equation}
  \rp{jk}{\ell}
    = \frac{\displaystyle\sum_{i} \tilde{X}_{ji}\,\tilde{X}_{ki}}
           {\sqrt{\displaystyle\sum_{i}\tilde{X}_{ji}^2}
            \sqrt{\displaystyle\sum_{i}\tilde{X}_{ki}^2}}
  \label{eq:parcial_def}
\end{equation}

Existe además una fórmula equivalente en términos de correlaciones simples que
evita calcular residuos explícitamente:

\begin{equation}
  \boxed{
    \rp{jk}{\ell}
      = \frac{r_{jk} - r_{j\ell}\,r_{k\ell}}
             {\sqrt{(1 - r_{j\ell}^2)(1 - r_{k\ell}^2)}}
  }
  \label{eq:parcial_formula}
\end{equation}

\noindent Esta fórmula será el instrumento central del ejercicio.

\subsection{Coeficiente de determinación $R^2$}

En la regresión de $X_j$ sobre un conjunto de regresores $\mathbf{Z}$:
\begin{equation}
  X_{ji} = \gamma_0 + \gamma_1 Z_{1i} + \cdots + \gamma_p Z_{pi} + v_i
  \label{eq:reg_gen}
\end{equation}

\noindent el \textbf{coeficiente de determinación} mide la fracción de la varianza
de $X_j$ explicada linealmente por $\mathbf{Z}$:

\begin{equation}
  R^2 = 1 - \frac{\text{SCR}}{\text{SCT}}
       = 1 - \frac{\sum \hat{v}_i^2}{\sum (X_{ji}-\bar{X}_j)^2}
  \label{eq:R2_def}
\end{equation}

\noindent donde SCR es la suma de cuadrados de los residuos y SCT la suma de
cuadrados total. Se tiene $0 \le R^2 \le 1$, con $R^2 = 1$ si y sólo si
el ajuste es perfecto ($\hat{v}_i = 0$ para todo $i$).

\subsection{Relación entre $R^2$ y correlaciones parciales}

Para la regresión de $X_1$ sobre $X_2$ y $X_3$ existe la siguiente identidad:

\begin{equation}
  \Rd{1}{23}
    = r_{12}^2 + r_{13.2}^2(1 - r_{12}^2)
    = r_{13}^2 + r_{12.3}^2(1 - r_{13}^2)
  \label{eq:R2_parciales}
\end{equation}

\noindent Es decir, el $R^2$ múltiple se puede descomponer como la contribución
marginal de cada regresor al conjunto.

% ══════════════════════════════════════════════════════════════════════════
\section{Análisis de la condición de multicolinealidad perfecta}

\subsection{Qué implica la ecuación~( \ref{eq:relacion} )}

La condición $\lambda_1 X_{1i} + \lambda_2 X_{2i} + \lambda_3 X_{3i} = 0$ con
$\lambda_j$ no todos nulos significa que las tres variables viven en un subespacio
de dimensión dos (o menor) del espacio $\mathbb{R}^n$. En particular, asumiendo
$\lambda_1 \neq 0$, podemos despejar:

\begin{equation}
  X_{1i} = -\frac{\lambda_2}{\lambda_1} X_{2i}
            -\frac{\lambda_3}{\lambda_1} X_{3i}
  \label{eq:despeje_X1}
\end{equation}

\noindent Es decir, $X_1$ es una \emph{combinación lineal exacta} de $X_2$ y
$X_3$, sin término de error. Esto es la definición formal de
\textbf{multicolinealidad perfecta}.

\subsection{Implicación para las medias}

Sumando~(\ref{eq:relacion}) sobre los $n$ individuos y dividiendo por $n$:

\begin{equation}
  \lambda_1 \bar{X}_1 + \lambda_2 \bar{X}_2 + \lambda_3 \bar{X}_3 = 0
  \label{eq:medias}
\end{equation}

Restando~(\ref{eq:medias}) de~(\ref{eq:relacion}) observación a observación:

\begin{equation}
  \lambda_1 x_{1i} + \lambda_2 x_{2i} + \lambda_3 x_{3i} = 0
  \label{eq:desviaciones}
\end{equation}

\noindent donde $x_{ji} = X_{ji} - \bar{X}_j$. \textbf{La relación lineal exacta
se preserva en las desviaciones respecto a la media.} Esto es crucial porque las
sumas necesarias para calcular correlaciones y $R^2$ se expresan en desviaciones.

% ══════════════════════════════════════════════════════════════════════════
\section{Cálculo de los coeficientes de determinación}

\subsection{\texorpdfstring{$\Rd{1}{23} = 1$}{R2-1.23 = 1}}

Queremos regresar $X_1$ sobre $X_2$ y $X_3$. Pero por~(\ref{eq:despeje_X1}),
la relación exacta

\begin{equation}
  X_{1i}
    = \underbrace{-\frac{\lambda_2}{\lambda_1}}_{=\,\gamma_2} X_{2i}
    + \underbrace{-\frac{\lambda_3}{\lambda_1}}_{=\,\gamma_3} X_{3i}
    + \underbrace{0}_{=\,v_i}
  \label{eq:reg_exacta_1}
\end{equation}

\noindent ya es la regresión de $X_1$ sobre $X_2$ y $X_3$ con residuos idénticamente
nulos. Por la definición~(\ref{eq:R2_def}):

\begin{equation}
  \Rd{1}{23}
    = 1 - \frac{\sum \hat{v}_i^2}{\sum x_{1i}^2}
    = 1 - \frac{0}{\sum x_{1i}^2}
    = \boxed{1}
  \label{eq:R2_1_23}
\end{equation}

\noindent \textit{Interpretación}: $X_2$ y $X_3$ explican el $100\%$ de la
variación de $X_1$. No queda variación sin explicar.

\subsection{\texorpdfstring{$\Rd{2}{13} = 1$}{R2-2.13 = 1}}

Asumiendo $\lambda_2 \neq 0$, de~(\ref{eq:relacion}) se tiene:

\begin{equation}
  X_{2i} = -\frac{\lambda_1}{\lambda_2} X_{1i}
            -\frac{\lambda_3}{\lambda_2} X_{3i}
  \label{eq:despeje_X2}
\end{equation}

\noindent La regresión de $X_2$ sobre $X_1$ y $X_3$ tiene residuos nulos, por lo que:

\begin{equation}
  \Rd{2}{13}
    = 1 - \frac{0}{\sum x_{2i}^2}
    = \boxed{1}
  \label{eq:R2_2_13}
\end{equation}

\subsection{\texorpdfstring{$\Rd{3}{12} = 1$}{R2-3.12 = 1}}

Asumiendo $\lambda_3 \neq 0$:

\begin{equation}
  X_{3i} = -\frac{\lambda_1}{\lambda_3} X_{1i}
            -\frac{\lambda_2}{\lambda_3} X_{2i}
  \label{eq:despeje_X3}
\end{equation}

\noindent Por el mismo argumento:

\begin{equation}
  \Rd{3}{12}
    = 1 - \frac{0}{\sum x_{3i}^2}
    = \boxed{1}
  \label{eq:R2_3_12}
\end{equation}

\begin{clave}
  \textbf{Resultado general.} Bajo multicolinealidad perfecta entre $k$ variables,
  el $R^2$ de la regresión de cualquiera de ellas sobre las demás $k-1$ es
  exactamente igual a~1. No hay grado de libertad para el error porque la
  relación no admite perturbaciones.
\end{clave}

% ══════════════════════════════════════════════════════════════════════════
\section{Cálculo de las correlaciones parciales}

Usaremos la fórmula~(\ref{eq:parcial_formula}) junto con las sumas en
desviaciones derivadas de~(\ref{eq:desviaciones}).

\subsection{Expresiones para las correlaciones simples}

Las correlaciones simples entre cualquier par de variables se calculan como:

\begin{equation}
  r_{jk} = \frac{\sum x_{ji} x_{ki}}{\sqrt{\sum x_{ji}^2 \sum x_{ki}^2}}
  \label{eq:rjk}
\end{equation}

De~(\ref{eq:desviaciones}):
\begin{equation}
  x_{1i} = -\frac{\lambda_2}{\lambda_1} x_{2i} - \frac{\lambda_3}{\lambda_1} x_{3i}
  \label{eq:x1_desv}
\end{equation}

Introducimos la notación compacta: $a = -\lambda_2/\lambda_1$ y
$b = -\lambda_3/\lambda_1$, de modo que $x_{1i} = a\,x_{2i} + b\,x_{3i}$.
Asimismo definamos:

\begin{align}
  S_{22} &= \sum x_{2i}^2, \quad
  S_{33} = \sum x_{3i}^2, \quad
  S_{23} = \sum x_{2i} x_{3i}
  \label{eq:sumas}
\end{align}

Entonces:

\begin{align}
  S_{11} &= \sum x_{1i}^2
           = a^2 S_{22} + 2ab\, S_{23} + b^2 S_{33}
  \label{eq:S11}\\
  S_{12} &= \sum x_{1i} x_{2i}
           = a\,S_{22} + b\,S_{23}
  \label{eq:S12}\\
  S_{13} &= \sum x_{1i} x_{3i}
           = a\,S_{23} + b\,S_{33}
  \label{eq:S13}
\end{align}

Por lo tanto las correlaciones simples son:

\begin{align}
  r_{12} &= \frac{S_{12}}{\sqrt{S_{11}S_{22}}}
           = \frac{a\,S_{22} + b\,S_{23}}{\sqrt{(a^2 S_{22}+2ab S_{23}+b^2 S_{33})S_{22}}}
  \label{eq:r12_expr}\\[6pt]
  r_{13} &= \frac{S_{13}}{\sqrt{S_{11}S_{33}}}
           = \frac{a\,S_{23} + b\,S_{33}}{\sqrt{(a^2 S_{22}+2ab S_{23}+b^2 S_{33})S_{33}}}
  \label{eq:r13_expr}\\[6pt]
  r_{23} &= \frac{S_{23}}{\sqrt{S_{22}S_{33}}}
  \label{eq:r23}
\end{align}

\subsection{Cálculo de \texorpdfstring{$\rp{12}{3}$}{r12.3}}

Aplicamos~(\ref{eq:parcial_formula}) con $j=1$, $k=2$, $\ell=3$:

\begin{equation}
  \rp{12}{3}
    = \frac{r_{12} - r_{13}\,r_{23}}
           {\sqrt{(1-r_{13}^2)(1-r_{23}^2)}}
  \label{eq:r12_3_formula}
\end{equation}

Para evaluar esto utilizamos la relación $x_{1i} = a x_{2i} + b x_{3i}$ y un
resultado auxiliar fundamental:

\begin{lema}
  \label{lema1}
  Si $x_{1i} = a x_{2i} + b x_{3i}$ para todo $i$, entonces:
  \begin{equation}
    r_{12} - r_{13} r_{23}
      = \frac{a S_{22}(1 - r_{23}^2)}
             {\sqrt{S_{11} S_{22}}}
    \label{eq:numerador}
  \end{equation}
\end{lema}

\begin{proof}
Expandimos el numerador $r_{12} - r_{13} r_{23}$:
\begin{align*}
  r_{12} - r_{13} r_{23}
  &= \frac{S_{12}}{\sqrt{S_{11}S_{22}}}
   - \frac{S_{13}}{\sqrt{S_{11}S_{33}}} \cdot \frac{S_{23}}{\sqrt{S_{22}S_{33}}}\\[4pt]
  &= \frac{1}{\sqrt{S_{11}S_{22}}}
     \left(S_{12} - \frac{S_{13}S_{23}}{S_{33}}\right)\\[4pt]
  &= \frac{1}{\sqrt{S_{11}S_{22}}}
     \left[(aS_{22}+bS_{23}) - \frac{(aS_{23}+bS_{33})S_{23}}{S_{33}}\right]\\[4pt]
  &= \frac{1}{\sqrt{S_{11}S_{22}}}
     \left[aS_{22} + bS_{23} - aS_{23}\frac{S_{23}}{S_{33}} - bS_{23}\right]\\[4pt]
  &= \frac{1}{\sqrt{S_{11}S_{22}}}
     \left[aS_{22} - a\frac{S_{23}^2}{S_{33}}\right]\\[4pt]
  &= \frac{a\,S_{22}}{\sqrt{S_{11}S_{22}}}
     \left[1 - \frac{S_{23}^2}{S_{22}S_{33}}\right]
   = \frac{a\,S_{22}}{\sqrt{S_{11}S_{22}}}(1 - r_{23}^2)
\end{align*}
\end{proof}

Ahora evaluamos el denominador de~(\ref{eq:r12_3_formula}). Para $1-r_{13}^2$:

\begin{align}
  1 - r_{13}^2
  &= 1 - \frac{S_{13}^2}{S_{11}S_{33}}
   = \frac{S_{11}S_{33} - S_{13}^2}{S_{11}S_{33}}
  \label{eq:1mr13}
\end{align}

Calculamos $S_{11}S_{33} - S_{13}^2$:
\begin{align}
  S_{11}S_{33}
  &= (a^2 S_{22}+2ab S_{23}+b^2 S_{33})S_{33} \notag\\
  S_{13}^2
  &= (aS_{23}+bS_{33})^2 = a^2 S_{23}^2 + 2ab S_{23}S_{33} + b^2 S_{33}^2 \notag\\[4pt]
  S_{11}S_{33} - S_{13}^2
  &= a^2(S_{22}S_{33} - S_{23}^2)
   = a^2 S_{22}S_{33}(1 - r_{23}^2)
  \label{eq:det_parcial}
\end{align}

Por tanto:
\begin{equation}
  1 - r_{13}^2 = \frac{a^2 S_{22}S_{33}(1-r_{23}^2)}{S_{11}S_{33}}
               = \frac{a^2 S_{22}(1-r_{23}^2)}{S_{11}}
  \label{eq:1mr13_final}
\end{equation}

Análogamente para $1 - r_{23}^2$ (ya conocida). Entonces el denominador de~(\ref{eq:r12_3_formula}) es:

\begin{align}
  \sqrt{(1-r_{13}^2)(1-r_{23}^2)}
  &= \sqrt{\frac{a^2 S_{22}(1-r_{23}^2)}{S_{11}} \cdot (1-r_{23}^2)} \notag\\
  &= \frac{|a|\,\sqrt{S_{22}}}{{\sqrt{S_{11}}}} (1-r_{23}^2)
  \label{eq:denom}
\end{align}

Finalmente, sustituyendo~(\ref{eq:numerador}) y~(\ref{eq:denom})
en~(\ref{eq:r12_3_formula}):

\begin{equation}
  \rp{12}{3}
    = \frac{\dfrac{a\,S_{22}(1-r_{23}^2)}{\sqrt{S_{11}S_{22}}}}
           {\dfrac{|a|\,\sqrt{S_{22}}}{\sqrt{S_{11}}}(1-r_{23}^2)}
    = \frac{a\,S_{22}(1-r_{23}^2)}{\sqrt{S_{11}S_{22}}}
      \cdot
      \frac{\sqrt{S_{11}}}{|a|\,\sqrt{S_{22}}(1-r_{23}^2)}
    = \frac{a}{|a|}
    = \pm 1
  \label{eq:r12_3_resultado}
\end{equation}

\begin{clave}
  \[\rp{12}{3} = \pm 1\]
  El signo es $+1$ si $a = -\lambda_2/\lambda_1 > 0$, y $-1$ si $a < 0$.
  Es decir, el signo depende de si $\lambda_1$ y $\lambda_2$ tienen signos
  iguales u opuestos.
\end{clave}

\subsection{Cálculo de \texorpdfstring{$\rp{13}{2}$}{r13.2}}

Por simetría del argumento anterior, ahora usamos la representación
$x_{1i} = ax_{2i} + bx_{3i}$ y controlamos por $X_2$ en lugar de $X_3$.

\begin{equation}
  \rp{13}{2}
    = \frac{r_{13} - r_{12}\,r_{23}}
           {\sqrt{(1-r_{12}^2)(1-r_{23}^2)}}
  \label{eq:r13_2_formula}
\end{equation}

Calculamos el numerador:
\begin{align}
  r_{13} - r_{12} r_{23}
  &= \frac{S_{13}}{\sqrt{S_{11}S_{33}}}
   - \frac{S_{12}}{\sqrt{S_{11}S_{22}}} \cdot \frac{S_{23}}{\sqrt{S_{22}S_{33}}} \notag\\[4pt]
  &= \frac{1}{\sqrt{S_{11}S_{33}}}
     \left(S_{13} - \frac{S_{12}S_{23}}{S_{22}}\right) \notag\\[4pt]
  &= \frac{1}{\sqrt{S_{11}S_{33}}}
     \left[(aS_{23}+bS_{33}) - \frac{(aS_{22}+bS_{23})S_{23}}{S_{22}}\right] \notag\\[4pt]
  &= \frac{b\,S_{33}}{\sqrt{S_{11}S_{33}}}(1 - r_{23}^2)
  \label{eq:num_r13_2}
\end{align}

Análogamente al procedimiento anterior:
\begin{equation}
  1 - r_{12}^2
    = \frac{b^2 S_{33}(1-r_{23}^2)}{S_{11}}
  \label{eq:1mr12}
\end{equation}

Por tanto el denominador es:
\begin{equation}
  \sqrt{(1-r_{12}^2)(1-r_{23}^2)}
    = \frac{|b|\,\sqrt{S_{33}}}{\sqrt{S_{11}}}(1-r_{23}^2)
  \label{eq:denom_r13_2}
\end{equation}

Y la correlación parcial:
\begin{equation}
  \rp{13}{2}
    = \frac{b\,S_{33}(1-r_{23}^2)/\sqrt{S_{11}S_{33}}}
           {|b|\,\sqrt{S_{33}}(1-r_{23}^2)/\sqrt{S_{11}}}
    = \frac{b}{|b|}
    = \pm 1
  \label{eq:r13_2_resultado}
\end{equation}

\begin{clave}
  \[\rp{13}{2} = \pm 1\]
  El signo es $+1$ si $b = -\lambda_3/\lambda_1 > 0$, y $-1$ en caso contrario.
\end{clave}

\subsection{Cálculo de \texorpdfstring{$\rp{23}{1}$}{r23.1}}

Para esta correlación parcial controlamos por $X_1$. Por simetría, escribimos:

\begin{equation}
  x_{2i} = -\frac{\lambda_1}{\lambda_2} x_{1i}
            -\frac{\lambda_3}{\lambda_2} x_{3i}
           \equiv c\, x_{1i} + d\, x_{3i}
  \label{eq:x2_desv}
\end{equation}

\noindent con $c = -\lambda_1/\lambda_2$ y $d = -\lambda_3/\lambda_2$.

Aplicamos la fórmula:
\begin{equation}
  \rp{23}{1}
    = \frac{r_{23} - r_{12}\,r_{13}}
           {\sqrt{(1-r_{12}^2)(1-r_{13}^2)}}
  \label{eq:r23_1_formula}
\end{equation}

Siguiendo exactamente el mismo procedimiento algebraico:
\begin{align}
  r_{23} - r_{12}r_{13}
  &= \frac{d\,S_{33}}{\sqrt{S_{22}S_{33}}}(1-r_{13}^2) \cdot
     \frac{1}{(1-r_{13}^2)} \cdot (\ldots)
  \notag
\end{align}

El cálculo es análogo a los anteriores y lleva al resultado:

\begin{equation}
  \rp{23}{1}
    = \frac{d}{|d|}
    = \pm 1
  \label{eq:r23_1_resultado}
\end{equation}

\begin{clave}
  \[\rp{23}{1} = \pm 1\]
  El signo depende del signo de $d = -\lambda_3/\lambda_2$.
\end{clave}

% ══════════════════════════════════════════════════════════════════════════
\section{Interpretación geométrica}

La correlación parcial $\rp{jk}{\ell}$ es la correlación entre los residuos de
dos proyecciones ortogonales. Concretamente:

\begin{itemize}[leftmargin=1.6em, itemsep=4pt]
  \item Sea $\mathbf{M}_\ell = \mathbf{I} - \mathbf{P}_\ell$ el proyector ortogonal
    al espacio columna de $X_\ell$ (la matriz que obtiene residuos de regresar
    sobre $X_\ell$).

  \item $\rp{jk}{\ell} = \Corr(\mathbf{M}_\ell \mathbf{x}_j,\, \mathbf{M}_\ell \mathbf{x}_k)$

  \item Bajo la condición~(\ref{eq:relacion}), $\mathbf{x}_j$ es combinación
    lineal de $\mathbf{x}_k$ y $\mathbf{x}_\ell$. Al proyectar fuera de $\mathbf{x}_\ell$,
    lo que queda de $\mathbf{x}_j$ es proporcional a lo que queda de $\mathbf{x}_k$:
    \[
      \mathbf{M}_\ell \mathbf{x}_j \propto \mathbf{M}_\ell \mathbf{x}_k
    \]
    Dos vectores proporcionales tienen correlación $\pm 1$.
\end{itemize}

% ══════════════════════════════════════════════════════════════════════════
\section{Consecuencias para la estimación MCO}

\subsection{Singularidad de \texorpdfstring{$\mathbf{X'X}$}{X'X}}

Considere el modelo $\mathbf{y} = \mathbf{X}\boldsymbol{\beta} + \mathbf{u}$
donde $\mathbf{X} = [\mathbf{x}_1 \;\mathbf{x}_2\;\mathbf{x}_3]$. La condición~(\ref{eq:relacion})
implica que existe $\boldsymbol{\lambda} = (\lambda_1,\lambda_2,\lambda_3)' \neq \mathbf{0}$
tal que:

\begin{equation}
  \mathbf{X}\boldsymbol{\lambda} = \lambda_1\mathbf{x}_1 + \lambda_2\mathbf{x}_2 + \lambda_3\mathbf{x}_3
  = \mathbf{0}
  \label{eq:nulo}
\end{equation}

\noindent Por lo tanto $\boldsymbol{\lambda}$ pertenece al espacio nulo de
$\mathbf{X}$, lo que implica:

\begin{equation}
  \det(\mathbf{X'X}) = 0
  \label{eq:det_cero}
\end{equation}

La demostración: si $\mathbf{X}\boldsymbol{\lambda} = \mathbf{0}$, entonces
$\mathbf{X'X}\boldsymbol{\lambda} = \mathbf{X'}\mathbf{0} = \mathbf{0}$. Luego
$\boldsymbol{\lambda}$ es un vector propio de $\mathbf{X'X}$ asociado al valor
propio 0, lo cual implica $\det(\mathbf{X'X}) = 0$.

\subsection{Indeterminación de los estimadores}

Como $\mathbf{X'X}$ es singular, \textbf{no existe} $(\mathbf{X'X})^{-1}$ y el
estimador MCO usual

\begin{equation}
  \hat{\boldsymbol{\beta}} = (\mathbf{X'X})^{-1}\mathbf{X'y}
  \label{eq:beta_mco}
\end{equation}

\noindent \textbf{no está definido}. Hay infinitas combinaciones de
$(\hat\beta_1, \hat\beta_2, \hat\beta_3)$ que minimizan la suma de cuadrados
de los residuos.

\subsection{Varianzas infinitas}

Si pudiéramos en algún sentido hablar de varianzas, la fórmula para el caso
de dos regresores con correlación $r$ da:

\begin{equation}
  \text{Var}(\hat\beta_j) = \frac{\sigma^2}{S_{jj}(1 - R_j^2)}
  \label{eq:var_vif}
\end{equation}

\noindent Como hemos demostrado que $R_j^2 = 1$ para cada variable, el denominador
es $S_{jj}(1 - 1) = 0$, y la varianza diverge:

\begin{equation}
  \text{Var}(\hat\beta_j)
    = \frac{\sigma^2}{S_{jj} \cdot 0}
    \to \infty
  \label{eq:var_inf}
\end{equation}

\noindent Esto confirma que los coeficientes individuales son imposibles de estimar
con precisión finita.

% ══════════════════════════════════════════════════════════════════════════
\section{Resumen de resultados}

\begin{table}[ht]
\centering
\caption{Correlaciones parciales y $R^2$ bajo multicolinealidad perfecta}
\label{tab:resumen}
\small
\begin{tabular}{lcc}
\toprule
\textbf{Cantidad} & \textbf{Fórmula usada} & \textbf{Valor} \\
\midrule
$\rp{12}{3}$ & $\dfrac{r_{12}-r_{13}r_{23}}{\sqrt{(1-r_{13}^2)(1-r_{23}^2)}}$ & $\pm 1$ \\[10pt]
$\rp{13}{2}$ & $\dfrac{r_{13}-r_{12}r_{23}}{\sqrt{(1-r_{12}^2)(1-r_{23}^2)}}$ & $\pm 1$ \\[10pt]
$\rp{23}{1}$ & $\dfrac{r_{23}-r_{12}r_{13}}{\sqrt{(1-r_{12}^2)(1-r_{13}^2)}}$ & $\pm 1$ \\[10pt]
\midrule
$\Rd{1}{23}$ & $1 - \text{SCR}/\text{SCT}$ & $1$ \\[4pt]
$\Rd{2}{13}$ & $1 - \text{SCR}/\text{SCT}$ & $1$ \\[4pt]
$\Rd{3}{12}$ & $1 - \text{SCR}/\text{SCT}$ & $1$ \\
\midrule
$\text{VIF}_j$ & $1/(1-R_j^2)$ & $\infty$ \\
$\det(\mathbf{X'X})$ & --- & $0$ \\
\bottomrule
\end{tabular}
\end{table}

\begin{clave}
  \textbf{Conclusión.} La situación descrita en el ejercicio corresponde a
  \textbf{multicolinealidad perfecta} en su grado máximo posible. Sus
  manifestaciones son cuatro:

  \begin{enumerate}[leftmargin=1.4em, itemsep=4pt]
    \item \textbf{Todas} las correlaciones parciales entre pares de variables,
      controlando por la tercera, son iguales a $\pm 1$ en valor absoluto.

    \item \textbf{Todos} los $R^2$ de las regresiones auxiliares son iguales a 1:
      cada variable es perfectamente explicada por las demás dos.

    \item El VIF de cada variable es \textbf{infinito}: no existe información
      independiente en ninguna de ellas.

    \item El estimador MCO \textbf{no existe} (la matriz $\mathbf{X'X}$ es
      singular). Los errores estándar de los coeficientes individuales son
      infinitos. Solo son estimables las combinaciones lineales de parámetros
      que son ortogonales al vector de colinealidad $\boldsymbol{\lambda}$.
  \end{enumerate}
\end{clave}

% ══════════════════════════════════════════════════════════════════════════
\appendix

\section{Apéndice A: Demostración del Lema~\ref{lema1} mediante álgebra de Frisch-Waugh}

El teorema de Frisch-Waugh-Lovell (FWL) establece que, en la regresión de $\mathbf{y}$
sobre $[\mathbf{X}_1\;\mathbf{X}_2]$, el vector de coeficientes de $\mathbf{X}_2$
es idéntico al de la regresión de $\mathbf{M}_1\mathbf{y}$ sobre
$\mathbf{M}_1\mathbf{X}_2$, donde $\mathbf{M}_1 = \mathbf{I} - \mathbf{X}_1(\mathbf{X}_1'\mathbf{X}_1)^{-1}\mathbf{X}_1'$.

Aplicado a la correlación parcial: $\rp{12}{3}$ es la correlación entre
$\mathbf{M}_3\mathbf{x}_1$ y $\mathbf{M}_3\mathbf{x}_2$. Dado que
$\mathbf{x}_1 = a\mathbf{x}_2 + b\mathbf{x}_3$:

\begin{equation}
  \mathbf{M}_3\mathbf{x}_1
    = a\,\mathbf{M}_3\mathbf{x}_2 + b\,\mathbf{M}_3\mathbf{x}_3
    = a\,\mathbf{M}_3\mathbf{x}_2 + b\,\mathbf{0}
    = a\,\mathbf{M}_3\mathbf{x}_2
  \label{eq:fwl}
\end{equation}

\noindent porque $\mathbf{M}_3\mathbf{x}_3 = \mathbf{0}$ (proyectar sobre el
complemento del espacio de $\mathbf{x}_3$ aniquila a $\mathbf{x}_3$). Por tanto
$\mathbf{M}_3\mathbf{x}_1$ y $\mathbf{M}_3\mathbf{x}_2$ son proporcionales con
constante $a$, y su correlación es exactamente $\text{signo}(a) = \pm 1$.

\section{Apéndice B: Factor de inflación de varianza y número de condición}

\subsection{VIF a partir de $R^2$ auxiliar}

En la regresión múltiple general, la varianza del $j$-ésimo coeficiente MCO es:

\begin{equation}
  \text{Var}(\hat\beta_j) = \frac{\sigma^2}{S_{jj}}\cdot\frac{1}{1-R_j^2}
  \label{eq:var_vif_general}
\end{equation}

\noindent El factor $\text{VIF}_j = 1/(1-R_j^2)$ cuantifica el multiplicador de
la varianza respecto al caso ideal de ortogonalidad ($R_j^2 = 0$, VIF$=1$).
Bajo multicolinealidad perfecta, $R_j^2 = 1$ y VIF$_j = \infty$.

\subsection{Número de condición}

El \textbf{número de condición} de $\mathbf{X'X}$ (o de su versión estandarizada)
es:
\begin{equation}
  \kappa = \sqrt{\frac{\lambda_{\max}(\mathbf{X'X})}{\lambda_{\min}(\mathbf{X'X})}}
  \label{eq:kappa}
\end{equation}

\noindent Bajo multicolinealidad perfecta, $\lambda_{\min} = 0$, por lo que
$\kappa \to \infty$. En la práctica, $\kappa > 30$ ya indica colinealidad severa.

\section{Apéndice C: Generalización a $k$ variables}

\begin{prop}
  Si $k$ variables satisfacen $\sum_{j=1}^k \lambda_j X_{ji} = 0$ para todo $i$
  (con $\boldsymbol{\lambda} \neq \mathbf{0}$), entonces:
  \begin{enumerate}[label=(\roman*), itemsep=3pt]
    \item $R^2$ de la regresión de cualquier $X_j$ sobre las demás $k-1$
      variables es igual a~1.
    \item Toda correlación parcial entre dos variables, controlando por
      cualquier subconjunto de las restantes, tiene valor absoluto~1.
    \item La matriz $\mathbf{X'X}$ de dimensión $k\times k$ tiene rango
      a lo sumo $k-1$ (al menos un valor propio es 0).
    \item El estimador MCO $(\mathbf{X'X})^{-1}\mathbf{X'y}$ no existe.
  \end{enumerate}
\end{prop}

\vspace{18pt}
\noindent\rule{\linewidth}{0.4pt}

\noindent{\small\textit{Nota.} Los resultados de este ejercicio son puramente
algebraicos: no requieren supuestos distribucionales sobre los errores.
La multicolinealidad perfecta es una propiedad de los datos
($\mathbf{X}$), no del proceso generador de~$\mathbf{y}$.}

\end{document}
